\documentclass{article}
\usepackage[T1]{fontenc}
\usepackage{lmodern}
\usepackage{graphicx}
\usepackage{amsmath,amssymb}
\usepackage[a4paper,margin=2cm]{geometry}
\usepackage[absolute,overlay]{textpos}
\setlength{\TPHorizModule}{1mm}
\setlength{\TPVertModule}{1mm}
\usepackage[indonesian]{babel}
\usepackage{hyperref}
\setlength{\emergencystretch}{2em}
% unit: Halaman 24

\begin{document}

% === Halaman 24 (python/native) ===

24

Kriptanalisis Caesar Cipher

• Caesar cipher mudah dipecahkan dengan exhaustive key search  (brute force) karena 

jumlah kuncinya sangat sedikit (hanya ada 26 kunci).

• Coba lakukan dekripsi dengan berbagai nilai k dari 0 sampai 25, lalu periksa apakah 

hasil dekripsi merupakan kata atau kalimat yang bermakna. Jika ya, maka diduga k 

adalah kuncinya.

• Untuk memastikan k adalah kunci yang benar, maka cobakan k untuk potongan 

kriptogram lainnya.

\end{document}
